\section{Serialization}
\label{sec:serialization}
Serialization is the process of transforming data that is used by an application into a format that can be sent over a network or stored in a database or a file.

Deserialization, on the other hand, is the procedure that reads data from an external source and transforms it into a runtime object.

\subsection{Java serialization}
\label{ssec:java-serialization}
All the Alchemist data structures implements the Serializable interface~\cite{Serializ95:online}, meaning that it is possible to transform every object into a byte stream. The byte stream can be either saved to disk or transported to another system, where it can finally be deserialized to recreate the original object.

However, there are several problems with Java serialization:

\begin{itemize}
	\item \textbf{Absence of multi-platform support}: Serialized objects can only be deserialized in a JVM based environment. Other programming languages that do not translates into bytecode like Javascript ot Python cannot deserialize objects that have been serialized using the \textit{java.io.Serializable} interface. As it will be mentioned afterwards, this is a key point that the serialization strategy needs to have.
	\item \textbf{Security vulnerabilities}: Java serialization has been found to have various security vulnerabilities, such as allowing untrusted data to be deserialized, which can result in data breaches or other types of attacks. There’s no way to know what is being deserialized before the decoding operation. It is possible that an attacker serializes a malicious object and send them to your application. Once the readObject() method is called, the malicious objects have already been instantiated.~\cite{AvoidJav95:online}
	\item \textbf{Version incompatibility}: Serialized objects can become incompatible with newer versions of the class, which can cause errors when deserializing the object. The Java serialization format is not extensible, so it is not possible to add new fields to a class without breaking existing serialized objects.
	\item \textbf{Increased memory usage}: Serialized objects can take up more memory than their non-serialized counterparts, which can be a problem for systems with limited memory resources.
	\item \textbf{Some complex data structure cannot be serialized}: It is difficult to store certain types of data structures, for example lists with circular references, or even a simple Java Optional~\cite{Optional58:online}.
	\item \textbf{Not human-readable}: The serialized object is saved as a binary stream, which is a sequence of bytes that is not easily readable by humans.
\end{itemize}
\section{Third party library}
\label{sec:third-party-library}
There are several libraries available for serializing objects using a JVM based language. Some of the libraries that were taken into consideration to use for in this project are listed here:
\begin{itemize}
	\item \textbf{kotlinx.serialization}: A Kotlin library for serialization and deserialization, developed by the Kotlin team at JetBrains~\cite{Kotlinko43:online};
	\item \textbf{moshi}: A JSON library for Kotlin and Java, developed by Square. It uses a simple API and supports a wide range of types, including lists, maps, and nullable values~\cite{squaremo84:online};
	\item \textbf{gson}: A Java library for serializing and deserializing JSON, developed by Google. It provides
	a very simple and easy-to-use API~\cite{googlegs30:online};
	\item \textbf{jackson}: A very popular library for JSON Serialization and deserialization for both Kotlin and Java~\cite{FasterXM63:online}.
\end{itemize}

In this thesis, the \textit{kotlinx.serialization} library was chosen, primarily because of the multi-platform support that it offers. Some of the key benefits are:

\begin{itemize}
	\item \textbf{Code reuse}: Developers can share code across multiple platforms, reducing the amount of code duplication and making it easier to maintain and update the code base.
	\item \textbf{Consistency}: By sharing code across platforms, multi-platform libraries can help to ensure consistency in the behavior and user experience of an application across different platforms.
	\item \textbf{Cost-effective}: The costs associated with developing and maintaining separate code bases for different platforms are drastically reduced.
	\item \textbf{Reduced Development Time}: By reusing code, multi-platform libraries can help to reduce the time and effort required to develop an application for multiple platforms.
	\item \textbf{Familiarity}: APIs do not change upon changing platform, resulting in less effort required for unskilled developers to work on the project.
\end{itemize}

When it comes to serialize objects, the format must be chosen carefully, because each format has its own strengths and weaknesses, and the best choice of format depends on the specific requirements of the project. There are many different serialization formats available, but some of the most common ones include:
\begin{itemize}
	\item \textbf{JSON (JavaScript Object Notation)}: A lightweight, text-based format that is widely used for data interchange. It is easy to read and write, and is supported by many programming languages and libraries~\cite{JSON74:online}.
	\item \textbf{XML (eXtensible Markup Language)}: A markup language that is widely used for data representation. It is more verbose than JSON, but offers a number of features that are useful for representing complex data structures and for defining custom markup~\cite{Extensib95:online}.
	\item \textbf{YAML (YAML Ain't Markup Language)}: A human-readable data serialization format. It is designed to be easy to read and write. It is often used to write configuration files~\cite{WhatisYA36:online}.
	\item \textbf{Protocol Buffers (Google Protocol Buffers)}: A language and a platform-neutral data serialization format developed by Google. It is compact and efficient in tems of size and performances~\cite{Protocol60:online}.
	\item \textbf{Avro} : A data serialization format that is similar to Protocol Buffers and Thrift, but with a focus on data schemas and more compact binary encoding~\cite{WhatisAv78:online}.
	\item \textbf{BSON (binary JSON)}: BSON is a binary-encoded serialization of JSON-like documents. BSON is more space-efficient than JSON, and includes additional data types such as Date and Binary~\cite{Explaini59:online}.
\end{itemize}

In this thesis, the JSON format was chosen for the following motivations:
\begin{itemize}
	\item \textbf{Text based}: It is easy to read and debug. This can be especially useful when working with large or complex data sets like the Alchemist ones.
	\item \textbf{Wide support}: JSON has become the de facto standard for data interchange on the web.
	\item \textbf{Simple and flexible}: It is easy to add or remove fields as needed. This can be useful in situations where data needs to be shared between systems that have different requirements or where the data structure may change over time.
\end{itemize}
Nevertheless, it is possible to change from using JSON to using a different format at any time with very few changes, as \textit{kotlinx.serialization} is really versatile and can manipulates the format under the hoods.

\section{kotlinx.serialization examples}
\label{sec:kotlinx-serialization-examples}
Classes that can be serialized must be marked explicitly. Kotlin Serialization does not use reflections, so it is impossible to mistakenly deserialize a class that was not intended to be serializable. By including the \textit{@Serializable} annotation, a class is declared to be serializable.

\lstinputlisting[
language=Kotlin,
caption={Serializable Class example},
label={lst:serializable},
]{listings/Serializable.kt}

\noindent The Kotlin Serialization plugin is instructed to automatically create and connect a serializer for the class if it is marked with the \textit{@Serializable} annotation. The generated serializer class will inherit from the KSerializer interface. It is possible to retrieve the instance using the \textit{serializer()} function on the class's companion object.

Serializers are composed by three pieces:

\begin{itemize}
	\item A \textbf{serialize} function that implements a SerializationStrategy. It is given a value to serialize along with an instance of Encoder. It represents a value as a series of primitives using the \textit{encodeXxx} functions of Encoder. For each of the primitive types that serialization supports, there is an \textit{encodeXxx}.
	\item A \textbf{deserialize} function that implements a DeserializationStrategy. It takes a Decoder instance and returns a deserialized value. It makes use of the Decoder's \textit{decodeXxx} functions, which are an exact replica of Encoder's corresponding ones.
	\item A \textbf{descriptor} property that must accurately describe what the \textit{encodeXxx} and \textit{decodeXxx} functions do. This way a format implementation knows in advance which encoding and decoding techniques the \textit{encodeXxx} and \textit{decodeXxx} functions call.
\end{itemize}

It is possible to create a custom serializer and bind it to a class. This is done with the \textit{@Serializable} annotation by adding the \textit{with} property value.

\lstinputlisting[
language=Kotlin,
caption={Custome Serializer example},
label={lst:customSerializer},
]{listings/CustomSerializer.kt}

There is crucial problem considering the available code base.
The Alchemist Simulator is a JVM project, most of the data structures are written in Java or Kotlin and are meant to be used only inside a JVM. There are also lot of dependencies with libraries which does not support multiplatform at all.\newline

\noindent It is certainly still possible to treat Alchemist data structures like third-party classes and write a custom Serializer to make the \textit{@Serializable} annotation to work.
It is still important to mention that creating custom serializers for all the data structure is not ideal for the following reasons:
\begin{itemize}
	\item A deep understanding of the serialization process and the underlying data structure is required. This can be difficult to learn and master.
	\item Writing complex code to handle the details of serializing and deserializing the data is required. This can be time-consuming and error-prone.
	\item Custom serializers may need to be updated or maintained over time as the data structure or requirements change. This also results in a lot of time consumption and possible bugs/errors.
	\item Custom serializers may be difficult for other developers to understand and use, especially if they are not well-documented or are not designed with reuse in mind.
\end{itemize}

There's another problem in using custom serializer in this particular use case. As the original structures are written using a JVM based language, the deserialization process can only work inside a JVM. This results in a compatibility problem, as existent classes cannot be used in the current project in a Kotlin/JS context.
It would still be possible to navigate inside the Json Object created starting from the serialized string. This approach is also not desirable because accessing a JSON object can be much more difficult and tedious compared to accessing a structured class in a programming language.
JSON objects do not have a fixed structure or defined set of fields. This can be more cumbersome and error-prone, especially if the JSON object is large or complex, or if there is'nt familiarity with the syntax of the JSON library being used. All of these information leads to the approach that has been used.

\subsubsection{Surrogates}
\label{sssec:surrogates}
The following Java class is one of the simplest \textit{alchemist-api} class.

\lstinputlisting[
language=Java,
caption={Molecule class},
label={lst:molecule},
]{listings/Molecule.java}

In this case, the Molecule interface is defining a contract for a class to have a \textit{getName()} method that returns the name of the Molecule as a String. For this example, the inheritance of the Dependency interface can be ignored.

It is impossible to use the Molecule interface in a Kotlin/JS project, because Java is not supported. A surrogate class can be created to use the Molecule information in a context which is not JVM based.
A surrogate class is a class that is used to serialize and deserialize another class (commonly a third party class which cannot be accessed). It acts as a "surrogate" for the original class, handling the serialization and deserialization process on behalf of the original class.

Surrogate classes are often used when the original class is not directly serializable, either because:
\begin{itemize}
	\item It is not marked with a serialization annotation;
	\item It has complex internal state that cannot be easily serialized;
	\item Third-Party Dependencies that does not support serialization are involved.
\end{itemize}
By using a surrogate class, it is also possible to customize the serialization process and selectively include or exclude certain fields or properties from the serialized representation of the original class.

A valid surrogate class for the \textit{Molecule} class is the following:

\lstinputlisting[
language=Kotlin,
caption={MoleculeSurrogate class},
label={lst:moleculeSurrogate},
]{listings/MoleculeSurrogate.kt}

The process of serialization for this class is really easy.
\lstinputlisting[
language=Kotlin,
caption={MoleculeSurrogate serialization},
label={lst:serializingMoleculeSurrogate},
]{listings/SerializingMoleculeSurrogate.kt}

The remaining step is the creation of a mapping function that can translate the content from a Molecule to a MoleculeSurrogate:

This is a Kotlin extension function that extends the Molecule class with a new function called \textit{toMoleculeSurrogate()}. Extension functions allow to add new functions to a class without having to inherit from the class or use any type of design pattern such as decorator.

The function takes no arguments and returns an instance of the \textit{MoleculeSurrogate} class, which is initialized with the name property of the \textit{Molecule} instance that the function is being called on.

\lstinputlisting[
language=Kotlin,
caption={MoleculeSurrogate mapping function},
label={lst:toMoleculeSurrogate},
]{listings/ToMoleculeSurrogate.kt}

Using a surrogate class and its associated serialization and deserialization logic can be used on any platform, not just on the Java Virtual Machine (JVM).

These steps must be repeated for every Alchemist data structure that must be serialized. Because of composition, most extension functions will use other objects \textit{toXSurrogate()} to mutate their internal state into the corresponding surrogate classes.

\section{Polymorphism}
\label{sec:polymorphism}

kotlinx.serialization can also deal with polymorphic class hierarchies. In Kotlin, polymorphism can be open or closed.
\begin{itemize}
	\item ~\cite{kotlinxs84:online}.
\end{itemize}


Kotlin serialization handle open and closed polymorphism in different ways. In this project, the serialization of the \textit{EnvironmentSurrogate} class required two different strategy for the serialization and deserialization of the two generic types of the class. The type parameters that the class declares are both covariant, this means that if a subclass is a subtype of a superclass, then an instance of the class parameterized with the subclass is considered a subtype of the class parameterized with the superclass. The type parameters are:
\begin{itemize}
	\item \textbf{TS}: It represent the type of the concentration surrogate. It is constrained to a subtype of "Any", meaning that it can be any non nullable type in Kotlin.
	In the original \textit{Environment} data structure this type parameter is named \textbf{T} and represents the type of the concentration.
	\item \textbf{PS}: It represent the type of the position surrogate. It is constrained to a subtype of "PositionSurrogate" which means the type must be derived from that interface. In the original \textit{Environment} data structure this type parameter is named \textbf{P} and represents the type of the position.
\end{itemize}
